\documentclass[11pt]{article}
\usepackage{mathblatt}
\def\mitzone{1}
\begin{document}
\blattkopf{Trigonometrische Funktionen}{Gesamt}
\weit
\verzeichniszeile{\verz{zone}{Kennst du schon (Nr.~1–13)} \verztrenn \verz{e1}{1 Einheitskreis und Bogenmaß (Nr.~14–15)} \verztrenn \verz{e2}{2 Sinus- und Kosinusfunktion und ihre Merkmale (Nr.~16–17)} \verztrenn \verz{e3}{3 Parameter und Transformationen (Nr.~18–20)} \verztrenn \verz{e4}{4 Periodische Vorgänge modellieren (Nr.~21–22)} \verztrenn \verz{abhaken}{Das kann ich}}
% Zone „Das kennst du schon“ – aus bank/trigonometrische-funktionen/zone.jsonl (zusammenbau v0.1)
\einheitenkopf[zone][Kennst du schon]{Das kennst du schon}
\setcounter{aufgabe}{0}

%% TODO Titel: Ich-kann-Satz fehlt in der Bank (Platzhalter: Kettenname) – Nr. 1
%% TODO Anweisung: ein Satz über der ersten Teilaufgabe fehlt in der Bank – Nr. 1
\begin{aufgabe}{Sinus, Kosinus und Tangens als Seitenverhältnisse im rechtwinkligen Dreieck, Wert zu einem Winkel mit dem Taschenrechner im Grad-Modus}
\begin{teile}
\teil Gegenkathete 2 cm, Hypotenuse 5 cm – wie groß ist sin $\alpha$? sin $\alpha = $ \leerfeld
\teil sin 37° mit dem Taschenrechner im Grad-Modus – auf zwei Stellen? \leerfeld
\end{teile}
\end{aufgabe}

%% TODO Titel: Ich-kann-Satz fehlt in der Bank (Platzhalter: Kettenname) – Nr. 2
%% TODO Anweisung: ein Satz über der ersten Teilaufgabe fehlt in der Bank – Nr. 2
\begin{aufgabe}{Taschenrechner}
\begin{teile}
\teil Winkel in Grad – welchen Modus stellst du ein?\\ \kreuz{DEG}\\ \kreuz{RAD}
\teil sin $\alpha = 0{,}42$ – wie groß ist $\alpha$? Nutze die Umkehrtaste und runde auf eine Stelle. $\alpha \approx$ \leerfeld[°]
\end{teile}
\end{aufgabe}

%% TODO Titel: Ich-kann-Satz fehlt in der Bank (Platzhalter: Kettenname) – Nr. 3
%% TODO Anweisung: ein Satz über der ersten Teilaufgabe fehlt in der Bank – Nr. 3
\begin{aufgabe}{Koordinatensystem mit vier Quadranten, Punkte eintragen und ablesen, Achsen benennen und einteilen}
\begin{teile}
\teil Lies die Koordinaten von A ab.

\begin{ksys}[xmin=-4,xmax=4,ymin=-4,ymax=4,ablesen]\punkt{3}{2}{A}\end{ksys}
A(\leerfeld | \leerfeld)
\teil Lies die Koordinaten von C ab.

\begin{ksys}[xmin=-4,xmax=4,ymin=-4,ymax=4,xstep=0.5,ystep=0.5,ablesen]\punkt{-1.5}{-2.5}{C}\end{ksys}
C(\leerfeld | \leerfeld)
\end{teile}
\end{aufgabe}

%% TODO Titel: Ich-kann-Satz fehlt in der Bank (Platzhalter: Kettenname) – Nr. 4
%% TODO Anweisung: ein Satz über der ersten Teilaufgabe fehlt in der Bank – Nr. 4
\begin{aufgabe}{Kreis mit Mittelpunkt und Radius zeichnen, Umfang mit Pi berechnen, Mittelpunktswinkel und Kreisbogen als Anteil des Vollkreises}
\begin{teile}
\teil Radius 3 cm – wie lang ist der Umfang? Gib ihn mit $\pi$ und gerundet an. $u = $ \leerfeld[cm]
\teil Radius 4 cm, Mittelpunktswinkel 45° – wie lang ist der Kreisbogen? Runde auf zwei Stellen. $b \approx$ \leerfeld[cm]
\end{teile}
\end{aufgabe}

%% TODO Titel: Ich-kann-Satz fehlt in der Bank (Platzhalter: Kettenname) – Nr. 5
%% TODO Anweisung: ein Satz über der ersten Teilaufgabe fehlt in der Bank – Nr. 5
\begin{aufgabe}{Winkel messen, zeichnen und benennen, Vollwinkel, gestreckter und rechter Winkel, Winkel über neunzig Grad als stumpfe Winkel}
\begin{teile}
\teil Wie groß ist ein gestreckter Winkel? \leerfeld[°]
\teil Zeichne am Strahl einen Winkel von 135° ein.

\winkelstrahl
\end{teile}
\end{aufgabe}

%% TODO Titel: Ich-kann-Satz fehlt in der Bank (Platzhalter: Kettenname) – Nr. 6
%% TODO Anweisung: ein Satz über der ersten Teilaufgabe fehlt in der Bank – Nr. 6
\begin{aufgabe}{Bruchteile eines Vollkreises als Bruch und als Dezimalzahl (halb, viertel, drittel)}
\begin{teile}
\teil Ein Viertel des Vollkreises – wie viel Grad? \leerfeld[°]
\teil 72° – welcher Anteil des Vollkreises? Gib ihn als Bruch und als Dezimalzahl an. \leerfeld
\end{teile}
\end{aufgabe}

%% TODO Titel: Ich-kann-Satz fehlt in der Bank (Platzhalter: Kettenname) – Nr. 7
%% TODO Anweisung: ein Satz über der ersten Teilaufgabe fehlt in der Bank – Nr. 7
\begin{aufgabe}{Wertetabelle anlegen, Wertepaare als Punkte eintragen, den Graphen als durchgehende Linie zeichnen, Funktionswert und Argument ablesen, Punktprobe}
\begin{teile}
\teil $y = 2x$ – fülle die Wertetabelle aus.

\wertetabelle{x}{y}{0,1,2,3}
\teil $y = x^2 - 1$ – fülle die Wertetabelle aus.

\wertetabelle{x}{y}{-2,-1,0,1,2}
\end{teile}
\end{aufgabe}

%% TODO Titel: Ich-kann-Satz fehlt in der Bank (Platzhalter: Kettenname) – Nr. 8
%% TODO Anweisung: ein Satz über der ersten Teilaufgabe fehlt in der Bank – Nr. 8
\begin{aufgabe}{Achseneinteilung wählen, wenn beide Achsen verschiedene Einheiten und Schrittweiten haben}
\begin{teile}
\teil Rechtsachse: 0 bis 60 Minuten auf 6 cm – wie viele Minuten je Zentimeter? \leerfeld[min]
\teil Rechtsachse: 0 bis 2,5 Stunden auf 10 cm – wie viel je Zentimeter? \leerfeld[h]
\end{teile}
\end{aufgabe}

%% TODO Titel: Ich-kann-Satz fehlt in der Bank (Platzhalter: Kettenname) – Nr. 9
%% TODO Anweisung: ein Satz über der ersten Teilaufgabe fehlt in der Bank – Nr. 9
\begin{aufgabe}{Achsensymmetrisch und punktsymmetrisch als Wörter, Symmetrieachse und Symmetriezentrum am Graphen}
\begin{teile}
\teil $y = x^2$ – welche Symmetrie?\\ \kreuz{achsensymmetrisch zur Hochachse}\\ \kreuz{punktsymmetrisch zum Ursprung}
\teil $y = (x - 2)^2$ – wo liegt die Symmetrieachse? $x = $ \leerfeld
\end{teile}
\end{aufgabe}

%% TODO Titel: Ich-kann-Satz fehlt in der Bank (Platzhalter: Kettenname) – Nr. 10
%% TODO Anweisung: ein Satz über der ersten Teilaufgabe fehlt in der Bank – Nr. 10
\begin{aufgabe}{Parameter in einer Funktionsgleichung erkennen und ihre Wirkung auf den Graphen beschreiben (Streckung, Stauchung, Verschiebung), Graph zu Gleichung zuordnen}
\begin{teile}
\teil Welche Parabel ist schmaler?\\ \kreuz{$y = 3x^2$}\\ \kreuz{$y = x^2$}
\teil $y = -0{,}25x^2$ – wie wirkt der Faktor vor $x^2$? \leerfeld
\end{teile}
\end{aufgabe}

%% TODO Titel: Ich-kann-Satz fehlt in der Bank (Platzhalter: Kettenname) – Nr. 11
%% TODO Anweisung: ein Satz über der ersten Teilaufgabe fehlt in der Bank – Nr. 11
\begin{aufgabe}{Größtwert, Kleinstwert und Spannweite einer Tabelle entnehmen, Mittelwert bilden}
\begin{teile}
\teil Temperaturen 4, 9, 7, 11, 6 (in °C) – Größtwert und Kleinstwert? \leerfeld; \leerfeld
\teil −2, 5, 3, −1, 4 (in °C) – wie groß ist der Mittelwert? \leerfeld[°C]
\end{teile}
\end{aufgabe}

%% TODO Titel: Ich-kann-Satz fehlt in der Bank (Platzhalter: Kettenname) – Nr. 12
%% TODO Anweisung: ein Satz über der ersten Teilaufgabe fehlt in der Bank – Nr. 12
\begin{aufgabe}{Taschenrechner – Fehler finden}
\begin{teile}
\teil Tom rechnet 8 · sin 35° und erhält −3,43. Finde den Fehler und rechne richtig.
\end{teile}
\end{aufgabe}

%% TODO Titel: Ich-kann-Satz fehlt in der Bank (Platzhalter: Kettenname) – Nr. 13
%% TODO Anweisung: ein Satz über der ersten Teilaufgabe fehlt in der Bank – Nr. 13
\begin{aufgabe}{Taschenrechner – selbst rechnen}
\begin{teile}
\teil 6 · cos 48° – auf zwei Stellen? Prüfe vorher den Modus. \leerfeld
\end{teile}
\end{aufgabe}

\clearpage
% Einheit 1 von 4 – Einheitskreis und Bogenmaß (bank/trigonometrische-funktionen/e1.jsonl, zusammenbau v0.1)
%% TODO Zweigzeile Teil 1: was hier gelernt wird (Satz in Schülersprache aus dem Einheitstitel) – Einheit 1
\einheitenkopf[e1]{Einheit 1 von 4 · Einheitskreis und Bogenmaß}
\zweigzeile{ab Kl. 10 · keine P10-Aufgabe}
\setcounter{aufgabe}{13}

%% TODO Titel: Ich-kann-Satz fehlt in der Bank (Platzhalter: Kettenname) – Nr. 14
%% TODO Anweisung: ein Satz über der ersten Teilaufgabe fehlt in der Bank – Nr. 14
\begin{aufgabe}{Einheitskreis und Bogenmaß}
%% TODO \rechenplatz{n}: Schrittzahl des Grundfalls fehlt in der Bank (nur bei fester Schreibform, 2.3 b)
\begin{teile}
\teil Grad oder Bogenmaß? Die Winkelangabe lautet 2,4. Kreuze das Maß und den passenden Modus an.\\ \kreuz{Gradmaß}\kreuz{Bogenmaß}\\ \kreuz{DEG}\kreuz{RAD}
\teil Zeichne den Punkt P zum Winkel 40° auf dem Einheitskreis ein.

\begin{ksys}[xmin=-1.2,xmax=1.2,ymin=-1.2,ymax=1.2,xstep=0.2,ystep=0.2]\funktionab{sqrt(abs(1-\x^2))}{}{-1}{1}\funktionab{-sqrt(abs(1-\x^2))}{}{-1}{1}\end{ksys}
\teil Miss den Winkel $\alpha$ zum Punkt P.

\einheitskreis[ohne]{100}
$\alpha \approx$ \leerfeld[°]
\teil Lies sin $\alpha$ und cos $\alpha$ für den Punkt P ab (eine Stelle).

\begin{ksys}[xmin=-1.2,xmax=1.2,ymin=-1.2,ymax=1.2,xstep=0.2,ystep=0.2,ablesen]\funktionab{sqrt(abs(1-\x^2))}{}{-1}{1}\funktionab{-sqrt(abs(1-\x^2))}{}{-1}{1}\punkt{0.6}{0.8}{P}\end{ksys}
sin $\alpha \approx$ \leerfeld, cos $\alpha \approx$ \leerfeld
\teil P liegt ganz oben auf dem Einheitskreis – Winkel, Sinuswert und Kosinuswert? $\alpha = $ \leerfeld[°], sin $\alpha = $ \leerfeld, cos $\alpha = $ \leerfeld
\teil Am Einheitskreis wurde sin 50° ≈ 0,8 abgelesen. Prüfe mit dem Taschenrechner im Grad-Modus (zwei Stellen).

\einheitskreis{50}
\leerfeld
\teil 160° – welches Vorzeichen haben Sinuswert und Kosinuswert? sin: \leerfeld, cos: \leerfeld
\teil sin $\alpha = 0{,}6$ – welche beiden Winkel im ersten Vollkreis ab 0° passen? Runde auf eine Stelle. $\alpha_1 \approx$ \leerfeld[°], $\alpha_2 \approx$ \leerfeld[°]
\teil 60° ist ein Sechstel des Vollwinkels. Begründe mit dem Umfang des Einheitskreises, dass das Bogenmaß $\frac{\pi}{3}$ ist.
\teil 40° ins Bogenmaß – als Vielfaches von $\pi$ und auf zwei Stellen? \leerfeld
\teil $\frac{3}{4}\pi$ ins Gradmaß – wie viel Grad? \leerfeld[°]
\teil Welches Bogenmaß gehört zu 60°?\\ \kreuz{$\frac{\pi}{6}$}\kreuz{$\frac{\pi}{3}$}\kreuz{$\frac{\pi}{2}$}\kreuz{$\frac{2\pi}{3}$}
\teil Der Rechner zeigt für sin 1 den Wert 0,017. Welcher Modus ist eingestellt, und was zeigt er im anderen Modus (zwei Stellen)? \leerfeld
\teil Eine Diagonale e wird mit dem Sinussatz berechnet: e = 5,3 · sin 128° : sin 31° (in cm). Berechne e auf zwei Stellen und zeige am Einheitskreis, warum es sin 128° gibt, obwohl kein rechtwinkliges Dreieck einen Winkel von 128° hat. \hfill (P10 2015 OS) \\ $e \approx$ \leerfeld[cm]
\teil a) 105° – wie groß im Bogenmaß (zwei Stellen)? b) $\frac{4}{3}\pi$ – wie groß im Gradmaß? c) Erkläre am Einheitskreis, woher sin 105° kommt, und gib den Wert auf zwei Stellen an. a) \leerfeld\quad b) \leerfeld[°]\quad c) \leerfeld
\end{teile}
\end{aufgabe}

%% TODO Titel: Ich-kann-Satz fehlt in der Bank (Platzhalter: Kettenname) – Nr. 15
%% TODO Anweisung: ein Satz über der ersten Teilaufgabe fehlt in der Bank – Nr. 15
\begin{aufgabe}{Einheitskreis und Bogenmaß – Fehler finden, Begründen}
\begin{teile}
\teil Lena berechnet sin 25° und erhält −0,13. Finde den Fehler und gib den richtigen Wert an.
\teil Begründe am Einheitskreis, warum sin $\alpha$ nie größer als 1 sein kann.
\end{teile}
\end{aufgabe}

\clearpage
% Einheit 2 von 4 – Sinus- und Kosinusfunktion und ihre Merkmale (bank/trigonometrische-funktionen/e2.jsonl, zusammenbau v0.1)
%% TODO Zweigzeile Teil 1: was hier gelernt wird (Satz in Schülersprache aus dem Einheitstitel) – Einheit 2
\einheitenkopf[e2]{Einheit 2 von 4 · Sinus- und Kosinusfunktion und ihre Merkmale}
\zweigzeile{ab Kl. 10 · keine P10-Aufgabe}
\setcounter{aufgabe}{15}

%% TODO Titel: Ich-kann-Satz fehlt in der Bank (Platzhalter: Kettenname) – Nr. 16
%% TODO Anweisung: ein Satz über der ersten Teilaufgabe fehlt in der Bank – Nr. 16
\begin{aufgabe}{Sinus- und Kosinusfunktion}
%% TODO \rechenplatz{n}: Schrittzahl des Grundfalls fehlt in der Bank (nur bei fester Schreibform, 2.3 b)
\begin{teile}
\teil Wo ist die Welle zu Ende? Welcher Abschnitt ist genau eine ganze Wiederholung?\\ \kreuz{von A bis B}\\ \kreuz{von A bis C}

\begin{ksys}[xmin=0,xmax=720,ymin=-1.5,ymax=1.5,xstep=90,ystep=0.5,ablesen]\funktion{sin(\x)}{f}\punkt{0}{0}{A}\punkt{180}{0}{B}\punkt{360}{0}{C}\end{ksys}
\teil Fülle die Wertetabelle für y = sin(x) aus (zwei Stellen).

\wertetabelle{x in °}{y}{0,45,90,135,180,225,270,315,360}
\teil Trage die Punkte der Tabelle ein und zeichne die Welle y = sin(x).

\wertetabelle[0,0.71,1,0.71,0,-0.71,-1,-0.71,0]{x in °}{y}{0,45,90,135,180,225,270,315,360} \begin{ksys}[xmin=0,xmax=360,ymin=-1.2,ymax=1.2,xstep=45,ystep=0.2]\end{ksys}
\teil Rechtsachse für einen ganzen Vollkreis ab 0°, 9 cm Platz – wie viele Grad je Zentimeter? Welche Zahlen schreibst du an? \leerfeld[°] je cm
\teil Lies sin 135° am Graphen ab (eine Stelle).

\begin{ksys}[xmin=0,xmax=360,ymin=-1.2,ymax=1.2,xstep=45,ystep=0.2,ablesen]\funktion{sin(\x)}{f}\end{ksys}
\leerfeld
\teil Lies ab: Bei welchen Winkeln im ersten Vollkreis ist sin x ≈ 0,7?

\begin{ksys}[xmin=0,xmax=360,ymin=-1.2,ymax=1.2,xstep=45,ystep=0.2,ablesen]\funktion{sin(\x)}{f}\end{ksys}
$x_1 \approx$ \leerfeld[°], $x_2 \approx$ \leerfeld[°]
\teil Liegt P(50|0,77) auf dem Graphen von y = sin(x)? Rechne im Grad-Modus. \janein
\teil Welche Werte kann y = sin(x) annehmen? \leerfeld
\teil Gib alle Nullstellen von y = sin(x) zwischen 0° und 720° an. Welche Regelmäßigkeit siehst du? \leerfeld
\teil Gib Hoch- und Tiefpunkt von y = sin(x) im ersten Vollkreis ab 0° an. H(\leerfeld | \leerfeld), T(\leerfeld | \leerfeld)
\teil Nach welchem Winkel wiederholt sich y = sin(x) zum ersten Mal ganz? \leerfeld[°]
\teil y = sin(x) zwischen 100° und 250° – steigt oder fällt die Kurve?\\ \kreuz{steigt}\kreuz{fällt}
\teil Beschreibe die Symmetrie der Sinuskurve und belege sie mit einem Wertepaar.
\teil Um wie viel Grad muss man die Sinuskurve nach links schieben, damit sie auf der Kosinuskurve liegt?

\begin{ksys}[xmin=0,xmax=360,ymin=-1.2,ymax=1.2,xstep=45,ystep=0.2,ablesen]\funktion{sin(\x)}{f}\funktion{cos(\x)}{g}\end{ksys}
\leerfeld[°]
\teil Welche Gleichung gehört zum Graphen?\\ \kreuz{y = sin(x)}\\ \kreuz{y = cos(x)}

\begin{ksys}[xmin=0,xmax=360,ymin=-1.2,ymax=1.2,xstep=45,ystep=0.2,ablesen]\funktion{cos(\x)}{h}\end{ksys}
\teil Fülle die Wertetabelle für y = sin(x) aus (zwei Stellen), zeichne den Graphen und gib Wertebereich, Nullstellen, Hoch- und Tiefpunkt, kleinste Periode und Symmetrie an.

\wertetabelle{x in °}{y}{-180,-135,-90,-45,0,45,90,135,180} \begin{ksys}[xmin=-180,xmax=180,ymin=-1.2,ymax=1.2,xstep=45,ystep=0.2]\end{ksys}
\end{teile}
\end{aufgabe}

%% TODO Titel: Ich-kann-Satz fehlt in der Bank (Platzhalter: Kettenname) – Nr. 17
%% TODO Anweisung: ein Satz über der ersten Teilaufgabe fehlt in der Bank – Nr. 17
\begin{aufgabe}{Sinus- und Kosinusfunktion – Fehler finden, Begründen, Darstellungswechsel}
\begin{teile}
\teil Mehmet zeichnet die Sinuskurve so, dass sie im Hochpunkt spitz zuläuft wie ein Dach. Finde den Fehler.
\teil Begründe, warum sich die Sinuskurve nach einem Vollkreis wiederholt.
\teil Gehört die Wertetabelle zu f oder zu g?\\ \kreuz{f}\kreuz{g}

\wertetabelle[1,0,-1,0,1]{x in °}{y}{0,90,180,270,360} \begin{ksys}[xmin=0,xmax=360,ymin=-1.2,ymax=1.2,xstep=45,ystep=0.2,ablesen]\funktion{sin(\x)}{f}\funktion{cos(\x)}{g}\end{ksys}
\end{teile}
\end{aufgabe}

\clearpage
% Einheit 3 von 4 – Parameter und Transformationen (bank/trigonometrische-funktionen/e3.jsonl, zusammenbau v0.1)
%% TODO Zweigzeile Teil 1: was hier gelernt wird (Satz in Schülersprache aus dem Einheitstitel) – Einheit 3
\einheitenkopf[e3]{Einheit 3 von 4 · Parameter und Transformationen}
\zweigzeile{ab Kl. 10 · keine P10-Aufgabe}
\setcounter{aufgabe}{17}

%% TODO Titel: Ich-kann-Satz fehlt in der Bank (Platzhalter: Kettenname) – Nr. 18
%% TODO Anweisung: ein Satz über der ersten Teilaufgabe fehlt in der Bank – Nr. 18
\begin{aufgabe}{Parameter und Transformationen}
%% TODO \rechenplatz{n}: Schrittzahl des Grundfalls fehlt in der Bank (nur bei fester Schreibform, 2.3 b)
\begin{teile}
\teil Höher oder schmaler? Was hat sich bei g gegenüber f geändert?\\ \kreuz{die Höhe (a)}\kreuz{die Breite (b)}

\begin{ksys}[xmin=0,xmax=360,ymin=-2,ymax=2,xstep=45,ystep=0.5,ablesen]\funktion{sin(\x)}{f}\funktion{2*sin(\x)}{g}\end{ksys}
\teil y = 4 · sin(x) – wie groß ist die Amplitude? \leerfeld
\teil y = 5 · sin(x) – zwischen welchen Werten schwingt die Welle? \leerfeld
\teil Lies die Amplitude am Graphen ab.

\begin{ksys}[xmin=0,xmax=360,ymin=-3,ymax=3,xstep=45,ystep=0.5,ablesen]\funktion{2.5*sin(\x)}{f}\end{ksys}
$a = $ \leerfeld
\teil y = −2 · sin(x) – was ist anders als bei y = 2 · sin(x)? \leerfeld
\teil y = sin(3 · x) – wie lang ist die Periode? \leerfeld[°]
\teil Lies die Periode ab und berechne b in y = sin(b · x).

\begin{ksys}[xmin=0,xmax=360,ymin=-2,ymax=2,xstep=45,ystep=0.5,ablesen]\funktion{sin(4*\x)}{f}\end{ksys}
$b = $ \leerfeld
\teil Skizziere den Graphen von y = 2 · sin(3 · x) im ersten Vollkreis ab 0°.

\begin{ksys}[xmin=0,xmax=360,ymin=-3,ymax=3,xstep=30,ystep=1]\end{ksys}
\teil Stelle die Gleichung y = a · sin(b · x) zum Graphen auf: erst Amplitude, dann Periode, dann b.

\begin{ksys}[xmin=0,xmax=480,ymin=-5,ymax=5,xstep=60,ystep=1,ablesen]\funktion{4*sin(1.5*\x)}{f}\end{ksys}
$y = $ \leerfeld
\teil Welche Gleichung gehört zu g?\\ \kreuz{y = 2 · sin(x)}\\ \kreuz{y = sin(4 · x)}\\ \kreuz{y = 4 · sin(x)}

\begin{ksys}[xmin=0,xmax=360,ymin=-3,ymax=3,xstep=45,ystep=1,ablesen]\funktion{2*sin(\x)}{f}\funktion{sin(4*\x)}{g}\end{ksys}
\teil Der Faktor vor sin steigt von 1 auf 4. Beschreibe, was mit der Welle passiert.
\teil y = sin(x) + 2 – wohin ist die Welle verschoben, zwischen welchen Werten schwingt sie? \leerfeld
\teil Um wie viel Grad muss man y = sin(x) nach links schieben, damit y = cos(x) entsteht? Schreibe die Gleichung mit c. \leerfeld
\teil y = 2 · sin(4 · x) – erste Nullstelle nach 0° und erste Hochstelle? \leerfeld
\teil a) Stelle zum Graphen f die Gleichung y = a · sin(b · x) auf. b) Skizziere daneben den Graphen von y = 2 · sin(1,5 · x).

\begin{ksys}[xmin=0,xmax=360,ymin=-4,ymax=4,xstep=45,ystep=1,ablesen]\funktion{3.5*sin(3*\x)}{f}\end{ksys} \begin{ksys}[xmin=0,xmax=480,ymin=-3,ymax=3,xstep=60,ystep=1]\end{ksys}
\end{teile}
\end{aufgabe}

%% TODO Titel: Ich-kann-Satz fehlt in der Bank (Platzhalter: Kettenname) – Nr. 19
%% TODO Anweisung: ein Satz über der ersten Teilaufgabe fehlt in der Bank – Nr. 19
\begin{aufgabe}{zwei Graphen mit verschiedenen Parametern vergleichen}
\begin{teile}
\teil Vergleiche f und g im Bild: Was ist gleich, was ist verschieden?

\begin{ksys}[xmin=0,xmax=360,ymin=-3,ymax=3,xstep=45,ystep=1,ablesen]\funktion{2*sin(\x)}{f}\funktion{2*sin(3*\x)}{g}\end{ksys}
\end{teile}
\end{aufgabe}

%% TODO Titel: Ich-kann-Satz fehlt in der Bank (Platzhalter: Kettenname) – Nr. 20
%% TODO Anweisung: ein Satz über der ersten Teilaufgabe fehlt in der Bank – Nr. 20
\begin{aufgabe}{Parameter und Transformationen – Fehler finden, Begründen}
\begin{teile}
\teil Lisa schreibt zu y = sin(4 · x): Die Periode ist 1440°, denn b = 4 macht die Welle breiter. Finde den Fehler.
\teil Begründe, warum ein größeres b die Welle schmaler macht.
\end{teile}
\end{aufgabe}

\clearpage
% Einheit 4 von 4 – Periodische Vorgänge modellieren (bank/trigonometrische-funktionen/e4.jsonl, zusammenbau v0.1)
%% TODO Zweigzeile Teil 1: was hier gelernt wird (Satz in Schülersprache aus dem Einheitstitel) – Einheit 4
\einheitenkopf[e4]{Einheit 4 von 4 · Periodische Vorgänge modellieren}
\zweigzeile{ab Kl. 10 · keine P10-Aufgabe}
\setcounter{aufgabe}{20}

%% TODO Titel: Ich-kann-Satz fehlt in der Bank (Platzhalter: Kettenname) – Nr. 21
%% TODO Anweisung: ein Satz über der ersten Teilaufgabe fehlt in der Bank – Nr. 21
\begin{aufgabe}{Periodische Vorgänge modellieren}
%% TODO \rechenplatz{n}: Schrittzahl des Grundfalls fehlt in der Bank (nur bei fester Schreibform, 2.3 b)
\begin{teile}
\teil Wiederholt sich das? Die Höhe des Ventils am Fahrradreifen bei gleichmäßiger Fahrt.\\ \kreuz{wiederholt sich}\kreuz{wiederholt sich nicht}
\teil Der Ausschlag eines Uhrpendels nach links und rechts – periodisch oder nicht? Begründe in einem Satz.
\teil Ein Leuchtfeuer ist 2 s hell, dann 4 s dunkel, dann wieder hell. Wie lang ist die Periode? \leerfeld[s]
\teil Temperatur in einem Gewächshaus – Größtwert und Kleinstwert?

\sachtabelle{lcccccc}{Uhrzeit & 2 & 6 & 10 & 14 & 18 & 22}{Temperatur in °C & 15 & 13 & 17 & 25 & 27 & 21}
\leerfeld; \leerfeld
\teil Größtwert 27 °C, Kleinstwert 13 °C – Mittellinie und Amplitude? \leerfeld; \leerfeld
\teil Mittellinie 20 °C, Amplitude 7 °C, Periode 24 h – wie lautet die Gleichung y = a · sin(b · t) + d? $y = $ \leerfeld
\teil Eine Gondel fährt zwischen 2 m und 34 m Höhe, eine Runde dauert 8 min. Welche Gleichung passt?\\ \kreuz{y = 16 · sin(45 · t) + 18}\\ \kreuz{y = 18 · sin(45 · t) + 16}\\ \kreuz{y = 16 · sin(8 · t) + 18}
\teil Für die Temperatur gilt y = 7 · sin(15 · t) + 20 (t in Stunden nach 8 Uhr, y in °C). Berechne y für t = 4 und deute den Wert.
\teil y = 7 · sin(15 · t) + 20 (t in Stunden nach 8 Uhr) – wann ist es am wärmsten? \leerfeld
\teil Lies mit Hilfslinien ab: Wie hoch ist die Gondel nach 1 Minute?

\begin{ksys}[xmin=0,xmax=16,ymin=0,ymax=36,xstep=1,ystep=2,xlabel=t in min,ylabel=Höhe in m,ablesen]\funktion{16*sin(45*\x)+18}{f}\end{ksys}
\leerfeld[m]
\teil Das Modell y = 7 · sin(15 · t) + 20 beschreibt einen Sommertag im Gewächshaus. Beurteile, ob es für eine ganze Woche passt, und nenne eine Grenze.
\teil Fülle die Tabelle aus: Welcher Funktionstyp wiederholt sich, welcher nimmt je Schritt immer gleich viel zu?

\sachtabelle{lcccc}{ & Gerade & Parabel & Exponentialkurve & Sinuskurve}{wiederholt sich & \leerzelle & \leerzelle & \leerzelle & \leerzelle \\ gleiche Zunahme je Schritt & \leerzelle & \leerzelle & \leerzelle & \leerzelle}
\teil Ein Guthaben von 500 € wächst jedes Jahr um 3 \%. Welcher Graph passt? Kreuze an und begründe für die beiden anderen, warum sie nicht passen. \hfill (P10 2019 OS)\\ \kreuz{A}\kreuz{B}\kreuz{C}

\begin{ksys}[xmin=0,xmax=40,ymin=0,ymax=2000,xstep=10,ystep=500,xlabel=Jahre,ylabel=€,ablesen]\funktion{500+15*\x}{A}\end{ksys} \begin{ksys}[xmin=0,xmax=40,ymin=0,ymax=2000,xstep=10,ystep=500,xlabel=Jahre,ylabel=€,ablesen]\funktion{500*1.03^\x}{B}\end{ksys} \begin{ksys}[xmin=0,xmax=40,ymin=0,ymax=2000,xstep=10,ystep=500,xlabel=Jahre,ylabel=€,ablesen]\funktion{500*(1.03^\x-1)}{C}\end{ksys}
\teil Die Tageslänge in einer Stadt schwankt zwischen 8 h und 16,4 h, die Periode ist 365 Tage; t zählt die Tage ab dem 21. März. a) Welche Gleichung passt: (1) y = 4,2 · sin(0,99 · t) + 12,2, (2) y = 12,2 · sin(0,99 · t) + 4,2 oder (3) y = 4,2 · sin(365 · t) + 12,2? b) Welcher Graph, A oder B, passt? c) Berechne y für t = 100 und deute den Wert im Sachzusammenhang.

\begin{ksys}[xmin=0,xmax=360,ymin=0,ymax=22,xstep=60,ystep=2,xlabel=t in Tagen,ylabel=Stunden,ablesen]\funktion{4.2*sin(0.99*\x)+12.2}{A}\funktion{8.4*sin(0.99*\x)+12.2}{B}\end{ksys}
\end{teile}
\end{aufgabe}

%% TODO Titel: Ich-kann-Satz fehlt in der Bank (Platzhalter: Kettenname) – Nr. 22
%% TODO Anweisung: ein Satz über der ersten Teilaufgabe fehlt in der Bank – Nr. 22
\begin{aufgabe}{Periodische Vorgänge modellieren – Fehler finden, Begründen}
\begin{teile}
\teil Zum Graphen der Höhe einer Riesenradgondel über der Zeit sagt Emma: Der Graph zeigt den Weg der Gondel – sie fährt eine Welle. Finde den Fehler.
\teil Begründe, warum ein Guthaben mit festem Zinssatz keine Sinuskurve ergibt.
\end{teile}
\end{aufgabe}

% Abhakseite „Das kann ich“ – Zone nur im Gesamt (\mitzone)
%% TODO Ich-kann-Sätze: Platzhalter wie in den Aufgabendateien
\begin{abhakseite}
\ifdefined\mitzone
\abhakgruppe{Kennst du schon}
\abhak{1}{Sinus, Kosinus und Tangens als Seitenverhältnisse im rechtwinkligen Dreieck, Wert zu einem Winkel mit dem Taschenrechner im Grad-Modus}
\abhak{2}{Taschenrechner}
\abhak{3}{Koordinatensystem mit vier Quadranten, Punkte eintragen und ablesen, Achsen benennen und einteilen}
\abhak{4}{Kreis mit Mittelpunkt und Radius zeichnen, Umfang mit Pi berechnen, Mittelpunktswinkel und Kreisbogen als Anteil des Vollkreises}
\abhak{5}{Winkel messen, zeichnen und benennen, Vollwinkel, gestreckter und rechter Winkel, Winkel über neunzig Grad als stumpfe Winkel}
\abhak{6}{Bruchteile eines Vollkreises als Bruch und als Dezimalzahl (halb, viertel, drittel)}
\abhak{7}{Wertetabelle anlegen, Wertepaare als Punkte eintragen, den Graphen als durchgehende Linie zeichnen, Funktionswert und Argument ablesen, Punktprobe}
\abhak{8}{Achseneinteilung wählen, wenn beide Achsen verschiedene Einheiten und Schrittweiten haben}
\abhak{9}{Achsensymmetrisch und punktsymmetrisch als Wörter, Symmetrieachse und Symmetriezentrum am Graphen}
\abhak{10}{Parameter in einer Funktionsgleichung erkennen und ihre Wirkung auf den Graphen beschreiben (Streckung, Stauchung, Verschiebung), Graph zu Gleichung zuordnen}
\abhak{11}{Größtwert, Kleinstwert und Spannweite einer Tabelle entnehmen, Mittelwert bilden}
\abhak{12}{Taschenrechner – Fehler finden}
\abhak{13}{Taschenrechner – selbst rechnen}
\fi
\abhakgruppe{Einheit 1 · Einheitskreis und Bogenmaß}
\abhak{14}{Einheitskreis und Bogenmaß}
\abhak{15}{Einheitskreis und Bogenmaß – Fehler finden, Begründen}
\abhakgruppe{Einheit 2 · Sinus- und Kosinusfunktion und ihre Merkmale}
\abhak{16}{Sinus- und Kosinusfunktion}
\abhak{17}{Sinus- und Kosinusfunktion – Fehler finden, Begründen, Darstellungswechsel}
\abhakgruppe{Einheit 3 · Parameter und Transformationen}
\abhak{18}{Parameter und Transformationen}
\abhak{19}{zwei Graphen mit verschiedenen Parametern vergleichen}
\abhak{20}{Parameter und Transformationen – Fehler finden, Begründen}
\abhakgruppe{Einheit 4 · Periodische Vorgänge modellieren}
\abhak{21}{Periodische Vorgänge modellieren}
\abhak{22}{Periodische Vorgänge modellieren – Fehler finden, Begründen}
\end{abhakseite}

\end{document}
